\chapter{Usuarios piloto}\label{anexo:usuarios_piloto}

\section{Introducción}

Durante la fase final del desarrollo se llevó a cabo una validación de la aplicación con usuarios piloto reales. El objetivo fue obtener retroalimentación directa de personas que no habían participado en el desarrollo, para identificar problemas de usabilidad, ajustes visuales necesarios, nuevas ideas de funcionalidad y comportamientos inesperados de la interfaz en condiciones de uso real.

La distribución se realizó mediante \textbf{Expo Go}, que permite ejecutar la aplicación sin publicarla en las tiendas oficiales. Durante las primeras pruebas se compartió un enlace de Expo con un grupo de usuarios voluntarios de distintos perfiles y con diferentes dispositivos.

Posteriormente, Expo modificó el comportamiento de carga de proyectos en Expo Go: desde el 12 de mayo de 2026, las actualizaciones publicadas mediante EAS Update solo pueden abrirse si la cuenta autenticada es propietaria del proyecto o pertenece a la organización correspondiente~\cite{expoGoLoading2026}. Por ello, para pruebas abiertas posteriores y demostraciones con evaluadores externos, la estrategia se ajustó a Expo Go con túnel (\texttt{npx expo start --tunnel}), manteniendo el backend desplegado en Render como entorno de datos.

\section{Perfil de los usuarios}

Participaron un total de \textbf{11 usuarios piloto}, seleccionados entre conocidos y contactos de la autora, buscando diversidad en cuanto a:

\begin{itemize}
    \item \textbf{Dispositivo}: 7 usuarios con iPhone (iOS) y 3 usuarios con distintos modelos de Android.
    \item \textbf{Perfil tecnológico}: desde personas habituadas al uso de aplicaciones móviles hasta perfiles con menor experiencia digital.
    \item \textbf{Interés en el dominio}: personas con interés en el ocio y la cultura de Sevilla, perfil objetivo de la aplicación.
\end{itemize}

\section{Metodología de recogida de feedback}

El feedback se recogió a través de conversaciones directas y mensajes, sin un cuestionario estructurado. Se pidió a los 11 participantes que exploraran la aplicación libremente y que comunicaran cualquier problema, sugerencia o impresión. Se tomaron notas de los comentarios más recurrentes y se agruparon por categoría.

Para estructurar mínimamente la sesión, se indicó a cada usuario que completase los siguientes flujos principales antes de explorar libremente el resto de la aplicación:

\begin{enumerate}
    \item Registro e inicio de sesión.
    \item Consulta del listado de eventos y acceso al detalle de un evento.
    \item Creación de un evento propio (público o privado).
    \item Uso del mapa interactivo de eventos cercanos.
    \item Envío de un mensaje en el chat de un evento.
\end{enumerate}

La tabla~\ref{tab:piloto_flujos} recoge el número de usuarios que completaron cada flujo satisfactoriamente.

\begin{table}[H]
\centering
\caption{Completitud de flujos principales por los usuarios piloto (n = 11)}
\label{tab:piloto_flujos}
\begin{tabular}{|l|c|c|}
\hline
\textbf{Flujo} & \textbf{Completado} & \textbf{\%} \\
\hline
Registro e inicio de sesión & 11 / 11 & 100\% \\
\hline
Consulta de eventos y detalle & 11 / 11 & 100\% \\
\hline
Creación de un evento & 11 / 11 & 100\% \\
\hline
Uso del mapa interactivo & 11 / 11 & 100\% \\
\hline
Chat de evento & 11 / 11 & 100\% \\
\hline
\end{tabular}
\end{table}

Los 11 usuarios completaron los cinco flujos. Las incidencias detectadas durante la realización de estos flujos (principalmente en dispositivos Android) se recogen en la sección siguiente.

\section{Resultados e incidencias detectadas}

\subsection{Problemas de compatibilidad con Android}

La incidencia más significativa detectada durante las pruebas piloto fue la \textbf{incompatibilidad de varios elementos de la interfaz en dispositivos Android}, reportada por los 3 usuarios con dispositivo Android (100\% de los usuarios Android). Esta incidencia no había sido detectada durante el desarrollo, que se realizó íntegramente sobre iOS.

Los principales problemas observados en Android fueron:

\begin{itemize}
    \item \textbf{Mapa no visible} (3/3 usuarios Android): el componente de mapa (OpenStreetMap) no se renderizaba en la pantalla de mapa de eventos en ningún dispositivo Android, mostrando una pantalla en blanco. El problema estaba relacionado con la configuración del proveedor de tiles en la compilación de Expo Go para Android.

    \item \textbf{Descuadres de pantalla} (2/3 usuarios Android): varios elementos de la interfaz (modales, cabeceras, botones flotantes y disposición de formularios) presentaban desajustes de posicionamiento en Android, causados por diferencias en el comportamiento del \textit{safe area} y la gestión del teclado virtual respecto a iOS.
\end{itemize}

Estos problemas fueron corregidos en una iteración adicional de la fase de cierre del proyecto. Supusieron una desviación no prevista en la planificación, documentada en el capítulo~\ref{cap:planificación} (sección de desviaciones) y en el capítulo~\ref{cap:conclusiones} (lecciones aprendidas).

\subsection{Sugerencias de mejora visual y de interacción}

Los usuarios piloto aportaron numerosas sugerencias de mejora relacionadas con aspectos visuales y de comodidad en la interacción. Las más recurrentes fueron:

\begin{itemize}
    \item \textbf{Ajustes en el tamaño y contraste de algunos textos}: en determinadas pantallas, el texto informativo secundario resultaba difícil de leer, especialmente en modo oscuro. Se realizaron ajustes de contraste y tamaño en los componentes afectados.

    \item \textbf{Mejoras en el flujo de creación de eventos}: algunos usuarios encontraron el formulario de creación de eventos extenso. Se reorganizaron los campos para que los datos más importantes (título, categoría, fecha y ubicación) aparecieran en primer lugar.

    \item \textbf{Navegación entre secciones}: varios usuarios indicaron que no encontraban intuitivamente algunas secciones (como ``Mis eventos'' o ``Guardados''). Se revisó la estructura del menú lateral para mejorar la accesibilidad de estas pantallas.

    \item \textbf{Feedback visual en acciones}: en algunas acciones (apuntarse a un evento, guardar, enviar mensaje) no había suficiente retroalimentación visual inmediata. Se añadieron indicadores de carga y mensajes de confirmación en los casos más relevantes.
\end{itemize}

\subsection{Ideas de nuevas funcionalidades}

Los usuarios piloto también propusieron funcionalidades que consideraban interesantes y que no estaban disponibles en la versión distribuida. La más mencionada fue:

\begin{itemize}
    \item \textbf{Módulo avanzado de amigos} (mencionado por 4/11 usuarios): varios usuarios valoraron positivamente el módulo de amistad existente, que permite buscar usuarios, entrar en sus perfiles, seguir o dejar de seguir, enviar mensajes privados y ver los eventos a los que asisten. Como mejora futura, echaron en falta una gestión de contactos todavía más completa: distinguir amistades cercanas, recibir sugerencias de planes basadas en personas de confianza y poder enviar invitaciones directas a eventos. Esta funcionalidad ampliaría la dimensión social existente. Se ha incluido en la sección de trabajo futuro del capítulo de conclusiones.
\end{itemize}

\section{Valoración general de los usuarios}

La valoración global de los 11 usuarios piloto fue muy positiva, sin quejas funcionales. Los aspectos destacados de forma unánime o mayoritaria fueron:

\begin{itemize}
    \item \textbf{Intención de descarga} (11/11 usuarios): todos los usuarios indicaron que se descargarían la aplicación si estuviera disponible en las tiendas oficiales y que la utilizarían en su día a día para descubrir planes en Sevilla.

    \item \textbf{Estética e interfaz} (11/11 usuarios): la interfaz fue valorada muy positivamente de forma unánime, con mención expresa a la calidad visual, la coherencia de las pantallas, el modo oscuro y la comodidad de uso general.

    \item \textbf{Mapa interactivo de eventos} (mencionado por 11/11 usuarios): el mapa de eventos fue el aspecto más destacado por todos los usuarios, que lo señalaron como la funcionalidad que más les gustó y que consideran más diferenciadora respecto a otras aplicaciones similares.

    \item \textbf{Sistema de rutas de eventos} (mencionado por 11/11 usuarios): las rutas recomendadas y la posibilidad de crear rutas propias fueron también señaladas por todos los usuarios como una funcionalidad original que no han visto en otras aplicaciones de ocio, y que consideran especialmente útil para planificar una salida con varias actividades.

    \item \textbf{Módulo social como diferenciador} (mencionado por 8/11 usuarios): el chat de evento, los mensajes privados y el módulo de amistad fueron valorados como novedosos y útiles respecto a otras aplicaciones de eventos.
\end{itemize}

Ninguno de los 11 usuarios reportó quejas sobre el funcionamiento de la aplicación ni problemas durante la realización de los flujos principales. Las únicas incidencias detectadas fueron las de compatibilidad con Android descritas en la sección anterior, que afectaron a la visualización pero no impidieron completar ningún flujo.

\section{Conclusiones del proceso piloto}

El proceso de validación con usuarios piloto fue valioso en dos dimensiones. Por un lado, permitió detectar y corregir la incidencia crítica de compatibilidad con Android antes de la entrega final. Por otro, proporcionó retroalimentación cualitativa sobre usabilidad e ideas de nuevas funcionalidades que han sido incorporadas al trabajo futuro del proyecto.

La experiencia confirma que las pruebas con usuarios reales en dispositivos reales son un complemento imprescindible de las pruebas técnicas automatizadas: los problemas de compatibilidad multiplataforma y los ajustes de usabilidad difícilmente se detectan sin pasar la aplicación por manos de personas ajenas al desarrollo.
